\documentclass[10pt,a4paper]{amsart}
\usepackage[margin=19mm]{geometry}
\usepackage{amsmath,amssymb,mathtools}
\usepackage[scr=rsfs,cal=euler]{mathalfa}
\usepackage{xfrac}
\usepackage[unicode,hidelinks,bookmarks=false]{hyperref}
\newcommand{\ie}{i.e.\ }
\newcommand{\cf}{cf.\ }
\newcommand{\resp}{respectively\ }
\newcommand{\Subsection}{\subsection}
\providecommand{\nobreakdash}{\nobreak\hbox{-}}
\setlength{\emergencystretch}{2em}
\multlinegap0pt
\usepackage{bm}
\usepackage{bbm}
\usepackage{dsfont}
\usepackage{xspace}
\usepackage{cancel}
\usepackage{mathtools}
\usepackage{amsthm}
\makeatletter
\@ifundefined{proposition}{\newtheorem{proposition}{Proposition}}{}
\@ifundefined{lemma}{\newtheorem{lemma}[proposition]{Lemma}}{}
\makeatother
\newcommand{\R}{\mathbb{R}}
\title[Metric Geometry of Lebesgue, Wasserstein, and Gromov-Wassers]{Metric Geometry of Lebesgue, Wasserstein, and Gromov-Wasserstein Spaces: Submetries, Curvature, and Geodesics}
\author{Martin Bauer, Facundo M\'emoli, Tom Needham, Mao Nishino}
\date{Source: arXiv:2608.11680v1 (CC BY 4.0)}
\begin{document}
\maketitle

\noindent\emph{Source and licence.} Metric Geometry of Lebesgue, Wasserstein, and Gromov-Wasserstein Spaces: Submetries, Curvature, and Geodesics. Version arXiv:2608.11680v1 (CC BY 4.0), distributed under the Creative Commons Attribution 4.0 International licence, as recorded in the arXiv licence metadata for this exact version and shown on the abstract page https://arxiv.org/abs/2608.11680v1. Full rights notice: licenses/arxiv-2608.11680-CC-BY-4.0.txt.\par\smallskip

\noindent\textit{Editorial scope.} Counted unit: the Lemma whose source label is \texttt{lem:map\_approximation} in the article (source0.tex lines 798--800, source label \texttt{lem:map\_approximation}), delivered together with the article's own complete original proof of it (source0.tex lines 803--842, one \texttt{proof} environment of 40 lines). No mathematics is written, repaired or re-derived: statement and proof are the article's own text line for line. Required context retained uncounted: none. Every definition, display and label of those lines is retained unchanged. Omitted from this package and not redistributed: 1--797, 801--802, 843--3013. Nothing inside a retained excerpt is omitted or altered: every retained body is delivered exactly as the article's lines are written, so no omission receipt of any kind accompanies this package. Citations: the retained excerpts carry no citation command at all, so no bibliography entry of the article is transcribed and no reference is redistributed. Reference adaptation: none was needed and none was made; every reference inside the delivered unit resolves inside it, so no unresolved reference or citation is produced. Printed slips kept literal: no printed word, symbol, hypothesis or number of the counted statement or of its proof was altered. Layout and class: the article's own class is \texttt{article} with packages including graphicx, amsmath, amssymb, amsfonts, amsthm, comment, biblatex, xcolor, tikz-cd, geometry, mathrsfs, appendix, mathtools, quiver; the wrapper is the lane's pinned \texttt{amsart} editorial wrapper at 10pt on A4, which restates the theorem-like environments the retained text needs and adds the article's own macro definitions that the retained text needs (\texttt{\textbackslash R}), with extra wrapper packages (bm, bbm, dsfont, xspace, cancel, mathtools). The delivered numbering is the unit's own. Assets: no retained span contains a figure environment, a graphics-inclusion command, a PDF-page inclusion, a table or a TikZ picture; no image file is copied into this package, the package declares no asset dependency and no image file is redistributed with it. Licence: version arXiv:2608.11680v1 of this article is distributed under the Creative Commons Attribution 4.0 International licence, as recorded in the arXiv licence metadata for this exact version and shown on the abstract page https://arxiv.org/abs/2608.11680v1. A full rights notice is delivered at licenses/arxiv-2608.11680-CC-BY-4.0.txt. Characters: every retained excerpt is pure ASCII, so the delivered engine, XeLaTeX with its default Latin Modern text font, reports no missing character.
\begin{lemma}\label{lem:map_approximation}
Let $\pi \in \Pi(\mathscr{L},\mathscr{L})$. There exists a sequence of maps $\phi_m \in \mathrm{Aut}(\mathbf{I})$ such that $(\mathrm{id}_\mathbf{I} \times \phi_m)_\ast \mathscr{L}$ converges weakly to $\pi$. 
\end{lemma}

\begin{proof}
    Let $m$ be a fixed integer, and throughout the rest of the proof, let $[2^m] = \{1,\ldots,2^m\}$. First, partition $[0,1)$  into intervals $(\mathbf{I}_i)_{i \in [2^m]}$, with $\mathbf{I}_i = [(i-1)/2^m, i/2^m)$. For each pair $(i,j)$ with $i,j \in [2^m]$, set $\pi_{i,j} = \pi(\mathbf{I}_i \times \mathbf{I}_j)$.

    Next, we refine the partitions. For a fixed $i_0 \in [2^m]$, partition $\mathbf{I}_{i_0}$ into half-open intervals $(\mathbf{J}_{i_0,j})_{j \in [2^m]}$ as follows:
    \[
    \mathbf{J}_{i_0,j} = [a_{j-1},a_{j}), \mbox{ where } a_0 = \frac{(i_0-1)}{2^m}, \, a_j = a_0 + \sum_{k=1}^{j} \pi_{i_0,k} \mbox{ for } j \in [2^m].
    \]
    This produces a partition, since $\sum_{k=1}^{2^m} \pi_{i_0,k} = \mathscr{L}(\mathbf{I}_{i_0})$, and has the property that $\mathscr{L}(\mathbf{J}_{i_0,j}) = \pi_{i_0,j}$.
    Note that $\mathbf{J}_{i_0,j}$ can potentially be the empty set, which we consider as a degenerate half-open interval.
    Similarly, for fixed $j_0 \in [2^m]$, partition $\mathbf{I}_{j_0}$ into half-open intervals $(\mathbf{K}_{i,j_0})_{i \in [2^m]}$ such that $\mathscr{L}(\mathbf{K}_{i,j_0}) = \pi_{i,j_0}$. 

    Finally, define $\phi_m:\mathbf{I} \to \mathbf{I}$ as follows. For each pair $(i,j)$ such that $\pi_{i,j} > 0$, the (nonempty) intervals $\mathbf{J}_{i,j}$ and $\mathbf{K}_{i,j}$ defined above have the same Lebesgue measure  $\pi_{i,j}$, so there is a unique orientation-preserving translation taking  $\mathbf{J}_{i,j}$ onto $\mathbf{K}_{i,j}$, and we define $\phi_m|_{\mathbf{J}_{i,j}}$ to be this translation. Then define $\phi_m(1) = 1$. By construction, $\phi_m \in \mathrm{Aut}(\mathbf{I})$.

    It remains to show that $(\mathrm{id}_\mathbf{I} \times \phi_m)_\ast \mathscr{L} \xrightarrow{m \to \infty} \pi$, weakly. Let $f:\mathbf{I} \times \mathbf{I} \to \R$ be a Lipschitz test function, with Lipschitz constant $L$, with respect to (say) the $\ell^\infty$-product metric on $\mathbf{I} \times \mathbf{I}$ (that is, the metric $\left((q,r),(s,t)\right) \mapsto \max\{|q-s|,|r-t|\}$, which metrizes the standard topology). To prove the claim, it suffices to show that 
    \[
    \int_{\mathbf{I} \times \mathbf{I}} f(s,t) d(\mathrm{id}_\mathbf{I}\times \phi_m)_\ast \mathscr{L}(s,t) \xrightarrow{m \to \infty} \int_{\mathbf{I} \times \mathbf{I}} f(s,t) d\pi(s,t)
    \]
    (since convergence for Lipschitz test functions is sufficient to prove weak convergence on a compact domain). Letting $t_i$ denote the midpoint of interval $\mathbf{I}_i$, we will first approximate the integrals on either side of this expression by the quantity $\sum_{i,j=1}^{2^m} f(t_i,t_j)\pi_{i,j}$.
    We have 
    \begin{align}
        &\left|\int_{\mathbf{I} \times \mathbf{I}} f(s,t) d(\mathrm{id}_\mathbf{I} \times \phi_m)_\ast \mathscr{L}(s,t) - \sum_{i,j=1}^{2^m} f(t_i,t_j)\pi_{i,j} \right| \nonumber \\
        & \qquad \qquad = \left|\sum_{i,j=1}^{2^m} \int_{\mathbf{I}_i \times \mathbf{I}_j} f(s,t) d(\mathrm{id}_\mathbf{I} \times \phi_m)_\ast \mathscr{L}(s,t) - \sum_{i,j=1}^{2^m}  f(t_i,t_j) \int_{\mathbf{I}_i \times \mathbf{I}_j} d(\mathrm{id}_\mathbf{I} \times \phi_m)_\ast \mathscr{L}(s,t) \right| \label{eqn:1d_lemma_1}\\
        & \qquad \qquad \leq \sum_{i,j=1}^{2^m} \int_{\mathbf{I}_i \times \mathbf{I}_j} |f(s,t) - f(t_i,t_j)| d(\mathrm{id}_\mathbf{I} \times \phi_m)_\ast \mathscr{L}(s,t) \nonumber \\
        &\qquad \qquad \leq  \sum_{i,j=1}^{2^m} \int_{\mathbf{I}_i \times \mathbf{I}_j} L \cdot \max\{|s-t_i|,|t-t_j|\} d(\mathrm{id}_\mathbf{I} \times \phi_m)_\ast \mathscr{L}(s,t) \label{eqn:1d_lemma_2} \\
        &\qquad \qquad \leq \frac{L}{2^{m+1}}\sum_{i,j=1}^{2^m} \int_{\mathbf{I}_i \times \mathbf{I}_j}  d(\mathrm{id}_\mathbf{I} \times \phi_m)_\ast \mathscr{L}(s,t) = \frac{L}{2^{m+1}}, \label{eqn:1d_lemma_3}
    \end{align}
    where \eqref{eqn:1d_lemma_1} uses the fact that
    \begin{equation*}
    \pi_{i,j} = \mathscr{L}(\mathbf{J}_{i,j}) = \mathscr{L}\left(\left\{t \in \mathbf{I}_i \mid \phi_m(t) \in \mathbf{I}_j \right\} \right) = (\mathrm{id}_\mathbf{I}\times \phi_m)_\ast \mathscr{L}(\mathbf{I}_i \times \mathbf{I}_j),
    \end{equation*}
    \eqref{eqn:1d_lemma_2} uses the Lipschitz property of $f$, and \eqref{eqn:1d_lemma_3} uses the bound $|s-t_i| \leq 1/2^{m+1}$ for all $s \in \mathbf{I}_i$. By a similar (in fact, slightly more straightforward) argument, we also have 
    \[
    \left|\int_{\mathbf{I} \times \mathbf{I}} f(s,t) d\pi (s,t) - \sum_{i,j=1}^{2^m} f(t_i,t_j)\pi_{i,j} \right| \leq \frac{L}{2^{m+1}}.
    \]
    By the triangle inequality, it follows that 
    \[
    \left| \int_{\mathbf{I} \times \mathbf{I}} f(s,t) d(\mathrm{id}_\mathbf{I} \times \phi_m)_\ast \mathscr{L}(s,t) - \int_{\mathbf{I} \times \mathbf{I}} f(s,t) d\pi (s,t)\right| \leq \frac{L}{2^m} \xrightarrow{m \to \infty} 0,
    \]
    completing the proof.
\end{proof}

\end{document}
